% --- Archon physics formalization source begin ---
% archon:physics
% archon:covers PhyXMiniProblems/problem_phyx_mini_0356.lean
% archon:source-report reports/phyx_mini/problem_phyx_mini_0356.source.json

\chapter{Physics problem phyx\_mini\_0356}
\label{ch:PhyXMiniProblems_problem_phyx_mini_0356}

\paragraph{Problem source.}
\#\# Physical scenario

In a quasi-static process \$A \textbackslash{}rightarrow B\$ (see diagram) in which no heat is exchanged with the environment, the mean pressure \$\textbackslash{}bar\{p\}\$ of a certain amount of gas is found to change with its volume \$V\$ according to the relation \$\$\textbackslash{}bar\{p\} = \textbackslash{}alpha V\textasciicircum{}\{-\textbackslash{}frac\{5\}\{3\}\}\$\$ where \$\textbackslash{}alpha\$ is a constant.  All processes take the system from macrostate \$A\$ to macrostate \$B\$.

\#\# Question

Find the net heat absorbed by this system in the process: The volume is increased and heat is supplied to cause the pressure to decrease linearly with the volume.

\#\# Answer choices

- A: A: 7950J

- B: B: 8950J

- C: C: 5500J

- D: D: 16000J

\#\# Figure caption (auxiliary; use the image as primary evidence)

The image is a diagram of a thermodynamic process represented on a pressure-volume (P-V) graph. 

\#\#\# Axes:
- **Vertical Axis (\textbackslash{}( \textbackslash{}bar\{p\} \textbackslash{}))**: Labeled as pressure in units of \textbackslash{}( 10\textasciicircum{}6 \textbackslash{}) dynes/cm². It ranges from 1 to 32.
- **Horizontal Axis (V)**: Labeled as volume in units of \textbackslash{}( 10\textasciicircum{}3 \textbackslash{}) cm³. It ranges from 1 to 8.

\#\#\# Key Points and Lines:
- **Points**: A, A', B, B'
  - A is at the top-left corner, aligned with a pressure of 32.
  - A' is directly below A on the 1 \textbackslash{}( \textbackslash{}times 10\textasciicircum{}3 \textbackslash{}) cm³ line.
  - B is at the bottom-right on the 8 \textbackslash{}( \textbackslash{}times 10\textasciicircum{}3 \textbackslash{}) cm³ line.
  - B' is directly above B on the same pressure level as A, along the vertical.
  
- **Process Paths**:
  - **Path AB'**: Horizontal line indicating a process at constant volume (isobaric).
  - **Path BA'**: Vertical line indicating a process at constant volume.
  - **Curved Path**: Represents an adiabatic process with the equation \textbackslash{}( \textbackslash{}bar\{p\} = \textbackslash{}alpha V\textasciicircum{}\{-5/3\} \textbackslash{}).

\#\#\# Arrows:
- Arrows indicate direction along paths: 
  - **a**, **b**, and **c** denote segments of different processes.
  
\#\#\# Labels:
- The letters \textbackslash{}textbf\{a\}, \textbackslash{}textbf\{b\}, \textbackslash{}textbf\{c\} indicate different parts of the processes.
- The graph includes an equation for the adiabatic process: \textbackslash{}( \textbackslash{}bar\{p\} = \textbackslash{}alpha V\textasciicircum{}\{-5/3\} \textbackslash{}).

The diagram depicts a cycle involving isobaric, isochoric, and adiabatic processes in a thermodynamic context.

\#\# Dataset metadata

Laws of Thermodynamics; Physical Model Grounding Reasoning, Numerical Reasoning

\paragraph{Recorded answer/context.}
A: A: 7950J

\paragraph{Figure/image path.}
/root/proposal\_for\_physic/science-mango/phyx\_mini\_run/phyx\_data/test\_image/356.png

\paragraph{Formalization target.}
create a compiling Lean file with sorry bodies at `PhyXMiniProblems/problem\_phyx\_mini\_0356.lean`. The Lean declarations must preserve the physical quantities, dimensions or dimensional roles, figure labels, governing-law hypotheses, and final relation expressed by this problem.
Use Mathlib/Physlib names found through LeanExplore where available. If a physics API is missing, introduce faithful local abstractions rather than scalar placeholder aliases.

\begin{theorem}[Physics formalization target]
\label{thm:physics:phyx_mini_0356:target}
\lean{PhyXMiniProblems.ProblemPhyXMini0356.netHeatAbsorbedAlongLinearPressurePathInJoules}
\uses{def:physics:phyx-mini-0356:phyxminiproblems-problemphyxmini0356-energyinjoules, def:physics:phyx-mini-0356:phyxminiproblems-problemphyxmini0356-gaspvprocesssetup, def:physics:phyx-mini-0356:phyxminiproblems-problemphyxmini0356-matchesproblemstatementandprimaryfigure, def:physics:phyx-mini-0356:phyxminiproblems-problemphyxmini0356-hasphysicalpvparameters, def:physics:phyx-mini-0356:phyxminiproblems-problemphyxmini0356-satisfiesstatedpressurevolumepathlaws, def:physics:phyx-mini-0356:phyxminiproblems-problemphyxmini0356-satisfiesquasistaticboundaryworklaws, def:physics:phyx-mini-0356:phyxminiproblems-problemphyxmini0356-satisfiesadiabaticnoheatcondition, def:physics:phyx-mini-0356:phyxminiproblems-problemphyxmini0356-satisfiesclosedsystemfirstlaw}
The assigned autoformalize agent should translate this physics problem into Lean declarations in the covered file, with theorem and lemma proof bodies written as `by sorry`.
\end{theorem}
\begin{proof}
This is an autoformalization task, not a proof task. Produce statements that can later be proved by the physics prover without weakening the physical model.
\end{proof}

% --- Archon declaration topology begin ---
\section{Declaration topology}
\begin{definition}[Volume Quantity]
  \label{def:physics:phyx-mini-0356:phyxminiproblems-problemphyxmini0356-volumequantity}
  \lean{PhyXMiniProblems.ProblemPhyXMini0356.VolumeQuantity}
  A signed physical volume, with dimension L³
\end{definition}
\begin{proof}
  This is the declared model definition of Volume Quantity.
\end{proof}
\begin{definition}[Pressure Quantity]
  \label{def:physics:phyx-mini-0356:phyxminiproblems-problemphyxmini0356-pressurequantity}
  \lean{PhyXMiniProblems.ProblemPhyXMini0356.PressureQuantity}
  A physical pressure, with dimension M L⁻¹ T⁻²
\end{definition}
\begin{proof}
  This is the declared model definition of Pressure Quantity.
\end{proof}
\begin{definition}[Energy Quantity]
  \label{def:physics:phyx-mini-0356:phyxminiproblems-problemphyxmini0356-energyquantity}
  \lean{PhyXMiniProblems.ProblemPhyXMini0356.EnergyQuantity}
  A signed physical energy, used for heat, work, and internal energy
\end{definition}
\begin{proof}
  This is the declared model definition of Energy Quantity.
\end{proof}
\begin{definition}[volume In Cubic Meters]
  \label{def:physics:phyx-mini-0356:phyxminiproblems-problemphyxmini0356-volumeincubicmeters}
  \lean{PhyXMiniProblems.ProblemPhyXMini0356.volumeInCubicMeters}
  \uses{def:physics:phyx-mini-0356:phyxminiproblems-problemphyxmini0356-volumequantity}
  Read a physical volume in cubic metres
\end{definition}
\begin{proof}
  This is the declared model definition of volume In Cubic Meters.
\end{proof}
\begin{definition}[volume In Thousands Of Cubic Centimeters]
  \label{def:physics:phyx-mini-0356:phyxminiproblems-problemphyxmini0356-volumeinthousandsofcubiccentimeters}
  \lean{PhyXMiniProblems.ProblemPhyXMini0356.volumeInThousandsOfCubicCentimeters}
  \uses{def:physics:phyx-mini-0356:phyxminiproblems-problemphyxmini0356-volumequantity, def:physics:phyx-mini-0356:phyxminiproblems-problemphyxmini0356-volumeincubicmeters}
  Read a physical volume in the horizontal-axis unit 10³ cm³. One such unit is one litre, or 10⁻³ m³
\end{definition}
\begin{proof}
  This is the declared model definition of volume In Thousands Of Cubic Centimeters.
\end{proof}
\begin{definition}[pressure In Pascals]
  \label{def:physics:phyx-mini-0356:phyxminiproblems-problemphyxmini0356-pressureinpascals}
  \lean{PhyXMiniProblems.ProblemPhyXMini0356.pressureInPascals}
  \uses{def:physics:phyx-mini-0356:phyxminiproblems-problemphyxmini0356-pressurequantity}
  Read a physical pressure in pascals
\end{definition}
\begin{proof}
  This is the declared model definition of pressure In Pascals.
\end{proof}
\begin{definition}[pressure In Millions Of Dynes Per Square Centimeter]
  \label{def:physics:phyx-mini-0356:phyxminiproblems-problemphyxmini0356-pressureinmillionsofdynespersquarecentimeter}
  \lean{PhyXMiniProblems.ProblemPhyXMini0356.pressureInMillionsOfDynesPerSquareCentimeter}
  \uses{def:physics:phyx-mini-0356:phyxminiproblems-problemphyxmini0356-pressurequantity, def:physics:phyx-mini-0356:phyxminiproblems-problemphyxmini0356-pressureinpascals}
  Read a physical pressure in the vertical-axis unit 10⁶ dyn/cm². Since 1 dyn/cm² = 0.1 Pa, one axis unit is 10⁵ Pa
\end{definition}
\begin{proof}
  This is the declared model definition of pressure In Millions Of Dynes Per Square Centimeter.
\end{proof}
\begin{definition}[energy In Joules]
  \label{def:physics:phyx-mini-0356:phyxminiproblems-problemphyxmini0356-energyinjoules}
  \lean{PhyXMiniProblems.ProblemPhyXMini0356.energyInJoules}
  \uses{def:physics:phyx-mini-0356:phyxminiproblems-problemphyxmini0356-energyquantity}
  Read heat, work, or internal energy in SI joules
\end{definition}
\begin{proof}
  This is the declared model definition of energy In Joules.
\end{proof}
\begin{definition}[stated Adiabatic Exponent]
  \label{def:physics:phyx-mini-0356:phyxminiproblems-problemphyxmini0356-statedadiabaticexponent}
  \lean{PhyXMiniProblems.ProblemPhyXMini0356.statedAdiabaticExponent}
  The dimensionless exponent printed on the adiabatic curve
\end{definition}
\begin{proof}
  This is the declared model definition of stated Adiabatic Exponent.
\end{proof}
\begin{definition}[Macrostate Label]
  \label{def:physics:phyx-mini-0356:phyxminiproblems-problemphyxmini0356-macrostatelabel}
  \lean{PhyXMiniProblems.ProblemPhyXMini0356.MacrostateLabel}
  The four state labels printed in the primary pressure--volume diagram
\end{definition}
\begin{proof}
  This is the declared model definition of Macrostate Label.
\end{proof}
\begin{definition}[Process Path]
  \label{def:physics:phyx-mini-0356:phyxminiproblems-problemphyxmini0356-processpath}
  \lean{PhyXMiniProblems.ProblemPhyXMini0356.ProcessPath}
  The four directed routes from A to B shown in the diagram. Paths a and c are the two rectangular routes, while path b is the requested diagonal
\end{definition}
\begin{proof}
  This is the declared model definition of Process Path.
\end{proof}
\begin{definition}[Axis Quantity]
  \label{def:physics:phyx-mini-0356:phyxminiproblems-problemphyxmini0356-axisquantity}
  \lean{PhyXMiniProblems.ProblemPhyXMini0356.AxisQuantity}
  Which thermodynamic quantity is assigned to a plotted axis
\end{definition}
\begin{proof}
  This is the declared model definition of Axis Quantity.
\end{proof}
\begin{definition}[Pressure Axis Scale]
  \label{def:physics:phyx-mini-0356:phyxminiproblems-problemphyxmini0356-pressureaxisscale}
  \lean{PhyXMiniProblems.ProblemPhyXMini0356.PressureAxisScale}
  The scale printed beside the pressure axis
\end{definition}
\begin{proof}
  This is the declared model definition of Pressure Axis Scale.
\end{proof}
\begin{definition}[Volume Axis Scale]
  \label{def:physics:phyx-mini-0356:phyxminiproblems-problemphyxmini0356-volumeaxisscale}
  \lean{PhyXMiniProblems.ProblemPhyXMini0356.VolumeAxisScale}
  The scale printed beneath the volume axis
\end{definition}
\begin{proof}
  This is the declared model definition of Volume Axis Scale.
\end{proof}
\begin{definition}[Path Geometry]
  \label{def:physics:phyx-mini-0356:phyxminiproblems-problemphyxmini0356-pathgeometry}
  \lean{PhyXMiniProblems.ProblemPhyXMini0356.PathGeometry}
  Geometric form of a complete route from A to B
\end{definition}
\begin{proof}
  This is the declared model definition of Path Geometry.
\end{proof}
\begin{definition}[Path Thermal Description]
  \label{def:physics:phyx-mini-0356:phyxminiproblems-problemphyxmini0356-paththermaldescription}
  \lean{PhyXMiniProblems.ProblemPhyXMini0356.PathThermalDescription}
  Thermal characterization stated for the two physically relevant paths
\end{definition}
\begin{proof}
  This is the declared model definition of Path Thermal Description.
\end{proof}
\begin{definition}[Pressure Volume Figure]
  \label{def:physics:phyx-mini-0356:phyxminiproblems-problemphyxmini0356-pressurevolumefigure}
  \lean{PhyXMiniProblems.ProblemPhyXMini0356.PressureVolumeFigure}
  \uses{def:physics:phyx-mini-0356:phyxminiproblems-problemphyxmini0356-macrostatelabel, def:physics:phyx-mini-0356:phyxminiproblems-problemphyxmini0356-processpath, def:physics:phyx-mini-0356:phyxminiproblems-problemphyxmini0356-axisquantity, def:physics:phyx-mini-0356:phyxminiproblems-problemphyxmini0356-pressureaxisscale, def:physics:phyx-mini-0356:phyxminiproblems-problemphyxmini0356-volumeaxisscale, def:physics:phyx-mini-0356:phyxminiproblems-problemphyxmini0356-pathgeometry}
  Qualitative and geometric information visible in the primary bitmap. The physical coordinate values themselves live in GasPVProcessSetup
\end{definition}
\begin{proof}
  This is the declared model definition of Pressure Volume Figure.
\end{proof}
\begin{definition}[Gas PVProcess Setup]
  \label{def:physics:phyx-mini-0356:phyxminiproblems-problemphyxmini0356-gaspvprocesssetup}
  \lean{PhyXMiniProblems.ProblemPhyXMini0356.GasPVProcessSetup}
  \uses{def:physics:phyx-mini-0356:phyxminiproblems-problemphyxmini0356-volumequantity, def:physics:phyx-mini-0356:phyxminiproblems-problemphyxmini0356-pressurequantity, def:physics:phyx-mini-0356:phyxminiproblems-problemphyxmini0356-energyquantity, def:physics:phyx-mini-0356:phyxminiproblems-problemphyxmini0356-macrostatelabel, def:physics:phyx-mini-0356:phyxminiproblems-problemphyxmini0356-processpath, def:physics:phyx-mini-0356:phyxminiproblems-problemphyxmini0356-paththermaldescription, def:physics:phyx-mini-0356:phyxminiproblems-problemphyxmini0356-pressurevolumefigure}
  Independent physical quantities for the gas and its processes. meanPressureTraceInPascals path v is the scalar readout of the path's mean pressure at SI volume coordinate v in cubic metres. The coefficient α is also stored as an SI scalar because its unit is the compound unit Pa·m⁵ once the stated exponent 5/3 is imposed. Heat and work remain independent dimensionful fields until the governing laws below relate them
\end{definition}
\begin{proof}
  This is the declared model definition of Gas PVProcess Setup.
\end{proof}
\begin{definition}[Between Endpoint Volumes]
  \label{def:physics:phyx-mini-0356:phyxminiproblems-problemphyxmini0356-betweenendpointvolumes}
  \lean{PhyXMiniProblems.ProblemPhyXMini0356.BetweenEndpointVolumes}
  \uses{def:physics:phyx-mini-0356:phyxminiproblems-problemphyxmini0356-volumeincubicmeters, def:physics:phyx-mini-0356:phyxminiproblems-problemphyxmini0356-gaspvprocesssetup}
  The interval of SI volume coordinates traversed from A to B
\end{definition}
\begin{proof}
  This is the declared model definition of Between Endpoint Volumes.
\end{proof}
\begin{definition}[Matches Problem Statement And Primary Figure]
  \label{def:physics:phyx-mini-0356:phyxminiproblems-problemphyxmini0356-matchesproblemstatementandprimaryfigure}
  \lean{PhyXMiniProblems.ProblemPhyXMini0356.MatchesProblemStatementAndPrimaryFigure}
  \uses{def:physics:phyx-mini-0356:phyxminiproblems-problemphyxmini0356-volumeinthousandsofcubiccentimeters, def:physics:phyx-mini-0356:phyxminiproblems-problemphyxmini0356-pressureinmillionsofdynespersquarecentimeter, def:physics:phyx-mini-0356:phyxminiproblems-problemphyxmini0356-statedadiabaticexponent, def:physics:phyx-mini-0356:phyxminiproblems-problemphyxmini0356-macrostatelabel, def:physics:phyx-mini-0356:phyxminiproblems-problemphyxmini0356-processpath, def:physics:phyx-mini-0356:phyxminiproblems-problemphyxmini0356-gaspvprocesssetup}
  Transcription of the primary image and prose. The endpoint data are A = (1 × 10³ cm³, 32 × 10⁶ dyn/cm²), B = (8 × 10³ cm³, 1 × 10⁶ dyn/cm²). The other two corners are A' and B'. Every arrowed route starts at A and ends at B; the requested route b is straight and heat-supplied, whereas the curved reference is adiabatic. No numerical heat, work, or internal-energy value occurs in this structure
\end{definition}
\begin{proof}
  This is the declared model definition of Matches Problem Statement And Primary Figure.
\end{proof}
\begin{definition}[Has Physical PVParameters]
  \label{def:physics:phyx-mini-0356:phyxminiproblems-problemphyxmini0356-hasphysicalpvparameters}
  \lean{PhyXMiniProblems.ProblemPhyXMini0356.HasPhysicalPVParameters}
  \uses{def:physics:phyx-mini-0356:phyxminiproblems-problemphyxmini0356-volumeincubicmeters, def:physics:phyx-mini-0356:phyxminiproblems-problemphyxmini0356-pressureinpascals, def:physics:phyx-mini-0356:phyxminiproblems-problemphyxmini0356-macrostatelabel, def:physics:phyx-mini-0356:phyxminiproblems-problemphyxmini0356-gaspvprocesssetup, def:physics:phyx-mini-0356:phyxminiproblems-problemphyxmini0356-betweenendpointvolumes}
  Positivity and nondegeneracy conditions for the physical process
\end{definition}
\begin{proof}
  This is the declared model definition of Has Physical PVParameters.
\end{proof}
\begin{definition}[Satisfies Stated Pressure Volume Path Laws]
  \label{def:physics:phyx-mini-0356:phyxminiproblems-problemphyxmini0356-satisfiesstatedpressurevolumepathlaws}
  \lean{PhyXMiniProblems.ProblemPhyXMini0356.SatisfiesStatedPressureVolumePathLaws}
  \uses{def:physics:phyx-mini-0356:phyxminiproblems-problemphyxmini0356-volumeincubicmeters, def:physics:phyx-mini-0356:phyxminiproblems-problemphyxmini0356-pressureinpascals, def:physics:phyx-mini-0356:phyxminiproblems-problemphyxmini0356-processpath, def:physics:phyx-mini-0356:phyxminiproblems-problemphyxmini0356-gaspvprocesssetup, def:physics:phyx-mini-0356:phyxminiproblems-problemphyxmini0356-betweenendpointvolumes}
  The path traces agree with their endpoint states. The curved reference obeys the displayed polytropic relation, and path b is the affine interpolation of the endpoint pressures as a function of volume. Thus “pressure decreases linearly with volume” is a physical path law, not a definition of the answer
\end{definition}
\begin{proof}
  This is the declared model definition of Satisfies Stated Pressure Volume Path Laws.
\end{proof}
\begin{definition}[Satisfies Quasistatic Boundary Work Laws]
  \label{def:physics:phyx-mini-0356:phyxminiproblems-problemphyxmini0356-satisfiesquasistaticboundaryworklaws}
  \lean{PhyXMiniProblems.ProblemPhyXMini0356.SatisfiesQuasistaticBoundaryWorkLaws}
  \uses{def:physics:phyx-mini-0356:phyxminiproblems-problemphyxmini0356-volumeincubicmeters, def:physics:phyx-mini-0356:phyxminiproblems-problemphyxmini0356-pressureinpascals, def:physics:phyx-mini-0356:phyxminiproblems-problemphyxmini0356-energyinjoules, def:physics:phyx-mini-0356:phyxminiproblems-problemphyxmini0356-processpath, def:physics:phyx-mini-0356:phyxminiproblems-problemphyxmini0356-gaspvprocesssetup}
  Quasistatic boundary-work laws, with work positive when done by the gas. For a straight path, work is the trapezoid area under the pressure trace. For a polytropic path p V\textasciicircum{}γ = constant, γ ≠ 1, integration gives W = (pᵢVᵢ - p\_fV\_f)/(γ - 1). These are general endpoint formulas and contain none of this problem's requested heat or derived numerical work values
\end{definition}
\begin{proof}
  This is the declared model definition of Satisfies Quasistatic Boundary Work Laws.
\end{proof}
\begin{definition}[Satisfies Adiabatic No Heat Condition]
  \label{def:physics:phyx-mini-0356:phyxminiproblems-problemphyxmini0356-satisfiesadiabaticnoheatcondition}
  \lean{PhyXMiniProblems.ProblemPhyXMini0356.SatisfiesAdiabaticNoHeatCondition}
  \uses{def:physics:phyx-mini-0356:phyxminiproblems-problemphyxmini0356-energyinjoules, def:physics:phyx-mini-0356:phyxminiproblems-problemphyxmini0356-processpath, def:physics:phyx-mini-0356:phyxminiproblems-problemphyxmini0356-gaspvprocesssetup}
  An adiabatic path exchanges no heat. The statement is quantified over any path carrying that thermal description and does not constrain the heat on the requested heat-supplied path b
\end{definition}
\begin{proof}
  This is the declared model definition of Satisfies Adiabatic No Heat Condition.
\end{proof}
\begin{definition}[Satisfies Closed System First Law]
  \label{def:physics:phyx-mini-0356:phyxminiproblems-problemphyxmini0356-satisfiesclosedsystemfirstlaw}
  \lean{PhyXMiniProblems.ProblemPhyXMini0356.SatisfiesClosedSystemFirstLaw}
  \uses{def:physics:phyx-mini-0356:phyxminiproblems-problemphyxmini0356-energyinjoules, def:physics:phyx-mini-0356:phyxminiproblems-problemphyxmini0356-processpath, def:physics:phyx-mini-0356:phyxminiproblems-problemphyxmini0356-gaspvprocesssetup}
  The closed-system first law on every route, with heat into the gas and work by the gas positive: Q = U\_finish - U\_start + W. Internal energy is stored by macrostate, so the same endpoint change is shared by every path from A to B without being separately postulated
\end{definition}
\begin{proof}
  This is the declared model definition of Satisfies Closed System First Law.
\end{proof}
\begin{definition}[Answer Choice]
  \label{def:physics:phyx-mini-0356:phyxminiproblems-problemphyxmini0356-answerchoice}
  \lean{PhyXMiniProblems.ProblemPhyXMini0356.AnswerChoice}
  The four heat values printed beside the answer-choice labels
\end{definition}
\begin{proof}
  This is the declared model definition of Answer Choice.
\end{proof}
\begin{definition}[displayed Heat In Joules]
  \label{def:physics:phyx-mini-0356:phyxminiproblems-problemphyxmini0356-displayedheatinjoules}
  \lean{PhyXMiniProblems.ProblemPhyXMini0356.displayedHeatInJoules}
  \uses{def:physics:phyx-mini-0356:phyxminiproblems-problemphyxmini0356-answerchoice}
  Displayed heat readout, in joules, for each answer choice
\end{definition}
\begin{proof}
  This is the declared model definition of displayed Heat In Joules.
\end{proof}
\begin{definition}[recorded Answer Choice]
  \label{def:physics:phyx-mini-0356:phyxminiproblems-problemphyxmini0356-recordedanswerchoice}
  \lean{PhyXMiniProblems.ProblemPhyXMini0356.recordedAnswerChoice}
  \uses{def:physics:phyx-mini-0356:phyxminiproblems-problemphyxmini0356-answerchoice}
  The answer label recorded in the supplied dataset metadata
\end{definition}
\begin{proof}
  This is the declared model definition of recorded Answer Choice.
\end{proof}
\begin{lemma}[adiabatic Reference Work And Internal Energy Change]
  \label{lem:physics:phyx-mini-0356:phyxminiproblems-problemphyxmini0356-adiabaticreferenceworkandinternalenergychange}
  \lean{PhyXMiniProblems.ProblemPhyXMini0356.adiabaticReferenceWorkAndInternalEnergyChange}
  \uses{def:physics:phyx-mini-0356:phyxminiproblems-problemphyxmini0356-energyinjoules, def:physics:phyx-mini-0356:phyxminiproblems-problemphyxmini0356-gaspvprocesssetup, def:physics:phyx-mini-0356:phyxminiproblems-problemphyxmini0356-matchesproblemstatementandprimaryfigure, def:physics:phyx-mini-0356:phyxminiproblems-problemphyxmini0356-hasphysicalpvparameters, def:physics:phyx-mini-0356:phyxminiproblems-problemphyxmini0356-satisfiesquasistaticboundaryworklaws, def:physics:phyx-mini-0356:phyxminiproblems-problemphyxmini0356-satisfiesadiabaticnoheatcondition, def:physics:phyx-mini-0356:phyxminiproblems-problemphyxmini0356-satisfiesclosedsystemfirstlaw}
  Along the adiabatic reference, the gas does 3600 J of work and absorbs no heat. Consequently the state-function change from A to B is -3600 J. Neither derived value is a premise of this lemma
\end{lemma}
\begin{proof}
  The conclusion follows from the governing relations and figure readouts named in the statement of adiabatic Reference Work And Internal Energy Change.
\end{proof}
\begin{lemma}[requested Straight Path Work In Joules]
  \label{lem:physics:phyx-mini-0356:phyxminiproblems-problemphyxmini0356-requestedstraightpathworkinjoules}
  \lean{PhyXMiniProblems.ProblemPhyXMini0356.requestedStraightPathWorkInJoules}
  \uses{def:physics:phyx-mini-0356:phyxminiproblems-problemphyxmini0356-energyinjoules, def:physics:phyx-mini-0356:phyxminiproblems-problemphyxmini0356-gaspvprocesssetup, def:physics:phyx-mini-0356:phyxminiproblems-problemphyxmini0356-matchesproblemstatementandprimaryfigure, def:physics:phyx-mini-0356:phyxminiproblems-problemphyxmini0356-hasphysicalpvparameters, def:physics:phyx-mini-0356:phyxminiproblems-problemphyxmini0356-satisfiesstatedpressurevolumepathlaws, def:physics:phyx-mini-0356:phyxminiproblems-problemphyxmini0356-satisfiesquasistaticboundaryworklaws}
  The area under the straight pressure trace is the trapezoid with endpoint pressures 32 × 10⁶ and 1 × 10⁶ dyn/cm² and volume change 7 × 10³ cm³, namely 11550 J
\end{lemma}
\begin{proof}
  The conclusion follows from the governing relations and figure readouts named in the statement of requested Straight Path Work In Joules.
\end{proof}
% --- Archon declaration topology end ---
% --- Archon physics formalization source end ---
